\chapter{Kỷ luật}
\markboth{Kỷ luật}{Discipline}

\section{Người không tự kéo mình dậy sẽ bị cuộc đời kéo đi.}
\begin{truyen}{Người không tự kéo mình dậy sẽ bị cuộc đời kéo đi.}{Wake up.}
\chuhoa{C}{âu thầy nói thẳng. Cả lớp im đi một chút, vì ai cũng hiểu: nếu mình không chủ động đứng dậy, hoàn cảnh sẽ đẩy mình theo cách đau hơn.}
\textit{(The teacher said it plainly. The class went quiet for a moment, because everyone understood: if you do not pull yourself up, life will drag you in harsher ways.)}
\end{truyen}

\begin{baihoc}
Chủ động với bản thân luôn nhẹ hơn bị hoàn cảnh ép phải thay đổi.
\textit{(Taking initiative with yourself is always easier than being forced by circumstances.)}
\end{baihoc}

\begin{ghinhoanh}
\textbf{Short English to remember:}

Lift yourself before life drags you.

\end{ghinhoanh}

\begin{tuhocnhanh}
1. Em đang chủ động hay đang bị kéo đi? \textit{Are you taking initiative or being dragged along?}

2. Điều gì em cần tự kéo mình dậy ngay? \textit{What do you need to pull yourself up from right now?}
\end{tuhocnhanh}
\ngancach

\section{Muốn đỡ khổ sau này, đừng chiều mình quá mức hôm nay.}
\begin{truyen}{Muốn đỡ khổ sau này, đừng chiều mình quá mức hôm nay.}{Wake up.}
\chuhoa{C}{ả lớp im khi nghe câu này. Nhiều cái chiều bản thân lúc nhỏ tưởng là dễ chịu, nhưng sau này lại hóa thành sự yếu đuối và thiếu bản lĩnh.}
\textit{(The class went quiet at this line. Many forms of self-indulgence feel pleasant now, but later turn into weakness and lack of resilience.)}
\end{truyen}

\begin{baihoc}
Dễ dãi hôm nay thường là món nợ ngày mai.
\textit{(Today's indulgence often becomes tomorrow's debt.)}
\end{baihoc}

\begin{ghinhoanh}
\textbf{Short English to remember:}

Too much comfort now creates pain later.

\end{ghinhoanh}

\begin{tuhocnhanh}
1. Em đang chiều mình ở điểm nào quá mức? \textit{Where are you being too indulgent with yourself?}

2. Nếu tiếp tục, điều gì sẽ xảy ra sau này? \textit{If you continue, what may happen later?}
\end{tuhocnhanh}
\ngancach

\section{Trễ một bài còn sửa được, trễ cả ý thức thì rất đắt.}
\begin{truyen}{Trễ một bài còn sửa được, trễ cả ý thức thì rất đắt.}{Wake up.}
\chuhoa{C}{âu nói làm vài em cúi xuống. Bài nộp trễ có thể xin thêm thời gian, nhưng khi một người để ý thức trách nhiệm xuống quá thấp, cái giá phải trả lớn hơn rất nhiều.}
\textit{(The sentence made a few students lower their heads. A late assignment may be fixed with more time, but when a person's sense of responsibility drops too low, the price becomes much greater.)}
\end{truyen}

\begin{baihoc}
Điều nguy hiểm không chỉ là việc chậm, mà là thói quen coi chậm trễ là bình thường.
\textit{(The danger is not only delay, but the habit of treating delay as normal.)}
\end{baihoc}

\begin{ghinhoanh}
\textbf{Short English to remember:}

Late work hurts less than late awareness.

\end{ghinhoanh}

\begin{tuhocnhanh}
1. Em có đang xem nhẹ sự chậm trễ của mình không? \textit{Are you taking your delays too lightly?}

2. Ý thức nào em cần kéo lên ngay? \textit{Which attitude do you need to raise immediately?}
\end{tuhocnhanh}
\ngancach

\section{Lười không mệt ngay, nhưng nó tính lãi rất cao.}
\begin{truyen}{Lười không mệt ngay, nhưng nó tính lãi rất cao.}{Wake up.}
\chuhoa{C}{ả lớp bật cười nhẹ vì câu nói giống lời cảnh báo tài chính. Nhưng ai cũng hiểu: lười hôm nay không đau liền, nên người ta chủ quan; đến lúc hậu quả đến thì đã cộng thêm rất nhiều “lãi”.}
\textit{(The class laughed lightly because the line sounded like financial advice. But everyone understood: laziness does not hurt immediately, so people become careless; by the time consequences arrive, heavy “interest” has already been added.)}
\end{truyen}

\begin{baihoc}
Hậu quả đến chậm không có nghĩa là hậu quả nhỏ.
\textit{(A delayed consequence is not a small consequence.)}
\end{baihoc}

\begin{ghinhoanh}
\textbf{Short English to remember:}

Laziness charges interest.

\end{ghinhoanh}

\begin{tuhocnhanh}
1. Em đang tích “lãi xấu” ở việc nào? \textit{In which area are you building bad interest?}

2. Em cần dừng ở đâu trước khi quá đắt? \textit{Where do you need to stop before it becomes too costly?}
\end{tuhocnhanh}
\ngancach

\section{Không ai cứu nổi người cứ thích trì hoãn bản thân.}
\begin{truyen}{Không ai cứu nổi người cứ thích trì hoãn bản thân.}{Wake up.}
\chuhoa{T}{hầy nói rất bình tĩnh: giúp một người là điều có thể, nhưng không ai có thể sống hộ phần quyết tâm của em. Nếu em cứ dời việc cần làm ra xa hơn, người khác khó mà kéo kịp.}
\textit{(The teacher spoke very calmly: helping someone is possible, but no one can live the part of your determination for you. If you keep pushing necessary work away, others can hardly pull you back in time.)}
\end{truyen}

\begin{baihoc}
Sự giúp đỡ chỉ có tác dụng khi em chịu hợp tác với chính tương lai của mình.
\textit{(Help only works when you cooperate with your own future.)}
\end{baihoc}

\begin{ghinhoanh}
\textbf{Short English to remember:}

No one can rescue repeated self-delay.

\end{ghinhoanh}

\begin{tuhocnhanh}
1. Em đang trì hoãn bản thân ở điều gì? \textit{Where are you delaying yourself?}

2. Ai đã giúp em mà em vẫn chưa chịu đổi? \textit{Who has tried to help you while you still refused to change?}
\end{tuhocnhanh}
\ngancach

\section{Em đang nuôi tương lai hay đang nuôi thói quen phá tương lai?}
\begin{truyen}{Em đang nuôi tương lai hay đang nuôi thói quen phá tương lai?}{Wake up.}
\chuhoa{C}{âu hỏi không cần thầy trả lời thay. Chỉ cần nhìn vào nếp sống mỗi ngày~--- ngủ nghỉ, dùng thời gian, học hành, chơi bời~--- là đủ thấy em đang nuôi điều gì lớn dần trong đời mình.}
\textit{(The teacher did not need to answer the question for them. Looking at daily habits alone—sleep, time use, study, entertainment—is enough to see what a person is feeding in life.)}
\end{truyen}

\begin{baihoc}
Mỗi thói quen là một khoản đầu tư; vấn đề là em đang đầu tư vào cái gì.
\textit{(Every habit is an investment; the question is what you are investing in.)}
\end{baihoc}

\begin{ghinhoanh}
\textbf{Short English to remember:}

Your habits are feeding your future.

\end{ghinhoanh}

\begin{tuhocnhanh}
1. Thói quen nào đang giúp tương lai em lớn lên? \textit{Which habit is helping your future grow?}

2. Thói quen nào đang phá nó? \textit{Which habit is damaging it?}
\end{tuhocnhanh}
\ngancach

\section{Cứ dễ dãi với mình mãi thì đời sẽ nghiêm khắc hộ.}
\begin{truyen}{Cứ dễ dãi với mình mãi thì đời sẽ nghiêm khắc hộ.}{Wake up.}
\chuhoa{C}{ả lớp lặng đi vì câu này rất thẳng. Khi ta bỏ qua quá nhiều giới hạn cần thiết, cuộc đời sẽ đặt ra những giới hạn đau hơn để bắt ta học điều mình lẽ ra phải học sớm.}
\textit{(The class grew quiet because the sentence was very direct. When we ignore too many necessary limits, life sets harsher limits so that we learn what we should have learned earlier.)}
\end{truyen}

\begin{baihoc}
Tự nghiêm với mình đúng lúc là cách tránh những cú dạy quá đau về sau.
\textit{(Being properly strict with yourself is how you avoid very painful lessons later.)}
\end{baihoc}

\begin{ghinhoanh}
\textbf{Short English to remember:}

Be strict with yourself before life has to be.

\end{ghinhoanh}

\begin{tuhocnhanh}
1. Em đang quá dễ dãi với mình ở điều gì? \textit{Where are you being too lenient with yourself?}

2. Nếu sửa sớm, em sẽ tránh được điều gì? \textit{What can you avoid by correcting it early?}
\end{tuhocnhanh}
\ngancach

\section{Đừng đợi bị bỏ lại rồi mới chịu chạy.}
\begin{truyen}{Đừng đợi bị bỏ lại rồi mới chịu chạy.}{Wake up.}
\chuhoa{T}{hầy không muốn học sinh chỉ biết quý nỗ lực khi đã quá muộn. Bắt đầu sớm có thể khó, nhưng luôn nhẹ hơn việc hốt hoảng chạy theo trong lo lắng và hụt hơi.}
\textit{(The teacher did not want students to value effort only when it was too late. Starting early may be hard, but it is always lighter than panicking and running after everything in fear.)}
\end{truyen}

\begin{baihoc}
Chủ động sớm là món quà em tặng cho chính mình của ngày mai.
\textit{(Acting early is a gift you give to your future self.)}
\end{baihoc}

\begin{ghinhoanh}
\textbf{Short English to remember:}

Run before the gap becomes too large.

\end{ghinhoanh}

\begin{tuhocnhanh}
1. Em đang chậm ở đâu? \textit{Where are you already falling behind?}

2. Em có thể bắt đầu đuổi kịp bằng việc gì ngay hôm nay? \textit{What can you start today to catch up?}
\end{tuhocnhanh}
\ngancach

\section{Ngại học hôm nay là đang đặt gánh nặng lên ngày mai.}
\begin{truyen}{Ngại học hôm nay là đang đặt gánh nặng lên ngày mai.}{Wake up.}
\chuhoa{C}{âu nói rất ngắn mà nặng. Mỗi lần trốn một phần việc hôm nay, em đang âm thầm giao thêm áp lực cho bản thân ngày mai, trong khi ngày mai vốn đã có việc riêng của nó.}
\textit{(The line was short but heavy. Every time you avoid part of today's work, you quietly hand extra pressure to tomorrow's self, even though tomorrow already has its own tasks.)}
\end{truyen}

\begin{baihoc}
Thương ngày mai thì phải làm tốt hơn hôm nay.
\textit{(If you care about tomorrow, you must do better today.)}
\end{baihoc}

\begin{ghinhoanh}
\textbf{Short English to remember:}

Tomorrow is already carrying enough.

\end{ghinhoanh}

\begin{tuhocnhanh}
1. Em có đang ném việc hôm nay cho ngày mai không? \textit{Are you throwing today's work at tomorrow?}

2. Việc nào em nên tự gánh ngay bây giờ? \textit{Which task should you carry right now?}
\end{tuhocnhanh}
\ngancach

\section{Sự thoải mái kéo dài thường là cái bẫy ngọt nhất.}
\begin{truyen}{Sự thoải mái kéo dài thường là cái bẫy ngọt nhất.}{Wake up.}
\chuhoa{C}{ả lớp hiểu vì sao câu này làm người lười dễ giật mình. Cái bẫy nguy nhất thường không đau ngay; nó mềm, ngọt, dễ chịu, nên mình cứ ở mãi trong đó cho tới khi nhận ra mình đã chậm đi quá nhiều.}
\textit{(The class understood why this line startles lazy students. The most dangerous trap often does not hurt immediately; it feels soft, sweet, and comfortable, so we stay in it until we realize we have fallen far behind.)}
\end{truyen}

\begin{baihoc}
Đừng nhầm cảm giác dễ chịu nhất thời với con đường tốt cho lâu dài.
\textit{(Do not confuse temporary comfort with a good long-term path.)}
\end{baihoc}

\begin{ghinhoanh}
\textbf{Short English to remember:}

Comfort can become a sweet trap.

\end{ghinhoanh}

\begin{tuhocnhanh}
1. Sự thoải mái nào đang giữ chân em? \textit{What comfort is holding you back?}

2. Em có dám bước ra khỏi nó không? \textit{Do you dare to step out of it?}
\end{tuhocnhanh}
\ngancach